\section{Ordered ordinary inertia on the new split component}
\label{sec:cubic-split-ordinary-points}

This section concerns the actual prime $P_*$ and embedding
$\Phi:U/P_*\hookrightarrow E[[T_1,\ldots,T_{12}]]$ of
AC-G54.  Here $\ell\ne2,3,269$ is split in $K$, the local
multiplier is one, and the seed is
$r_*=\operatorname{diag}(\xi\gamma,1,\xi\gamma^{-1})$,
where $\xi=\chi^j$ is unramified and
$\gamma=\psi\kappa_\ell^e$, $e\geq1$.
We prove the missing ordered ordinary condition on this
component; the fixed-type condition alone was not used as
a substitute for that assertion in AC-G54.

The published local input is Caraiani--Newton,
\emph{On the modularity of elliptic curves over imaginary
quadratic fields}, author version
\url{https://www.ma.imperial.ac.uk/~acaraian/papers/Modularity-IQF.pdf},
Definition~3.3.1 and Lemma~3.3.2, printed pp.~52--53.
The lemma states that a semistable representation over a
finite local coefficient algebra, of fixed regular Hodge
type and with semistable-ordinary residue representation,
is itself semistable-ordinary.  It is an all-Artin statement.
Its printed proof uses the canonical filtration by slopes
on the semistable period module and weak admissibility.
It is used as an established local theorem separately from
A.  We do not claim to reprove all of the underlying
$p$-adic comparison theory.  Its ordinary characters on
inertia are those specified in Definition~3.3.1.

\begin{lemma}[AC-G57-Artin: ordered inertia in the formal germ]
\label{AC-G57-Artin}
Let $C$ be a finite local $E$-algebra with residue field $E$.
Suppose $r_C:G_{\mathbb Q_\ell}\to\mathrm{GL}_3(C)$
is a fixed-type potentially crystalline deformation of
$r_*$, of the type in AC-G54.
Then $C^3$ has a full flag by free direct summands,
stable under $G_{\mathbb Q_\ell}$, whose ordered
inertia characters are
\[
                    \gamma,\quad1,\quad\gamma^{-1}.
                                                               \tag{G57.1}
\]
The flag is unique among flags with these ordered
characters on any open inertia subgroup killing $\psi$.
\end{lemma}
\begin{proof}
Choose a finite Galois extension $L/\mathbb Q_\ell$
over which $r_C$ is crystalline, and enlarge it to
kill $\psi$ and $\xi$.  Potential crystallinity over
the finite algebra has this meaning on its underlying
finite-dimensional coefficient representation; enlarging
$L$ preserves crystallinity.  Its restriction has the
parallel fixed Hodge type of
$\kappa_\ell^e\oplus1\oplus\kappa_\ell^{-e}$.
The reduction has its standard ordinary flag.
After enlarging $E$ to contain the embeddings of $L$
if necessary, Caraiani--Newton Lemma~3.3.2 applies.
It supplies an ordinary $G_L$-stable flag with
ordered characters on $I_L$
\[
                  \kappa_\ell^e,\quad1,\quad\kappa_\ell^{-e}.
                                                               \tag{G57.2}
\]
Indeed the Hodge type is parallel, so the product of
the local algebraic embedding characters in
Definition~3.3.1 is the corresponding power of the
cyclotomic character.  Equivalently this order is the
order of the standard ordinary flag of the residue
representation, to which the lemma applies.  In the
usual convention in which the cyclotomic character has
Hodge--Tate weight $-1$, its dominant weight is
$(e-2,-1,-e)$ at every embedding; it is dominant for
every $e\geq1$.

We give the descent and uniqueness argument explicitly.
Choose $s\in I_L$ such that
$d_0=\kappa_\ell(s)^e$, $d_1=1$, and
$d_2=\kappa_\ell(s)^{-e}$ are pairwise distinct.
Such an $s$ exists because the cyclotomic image on
every open inertia subgroup is infinite.
All $d_i-d_j$ are nonzero elements of $E$ and hence
units in $C$.  An upper triangular matrix with these
three diagonal entries has characteristic polynomial
$\prod_i(X-d_i)$; its Lagrange polynomials give three
orthogonal idempotents
\[
             e_i=\prod_{k\ne i}\frac{r_C(s)-d_k}{d_i-d_k}.
                                                               \tag{G57.3}
\]
They decompose $C^3$ into free rank-one eigenspaces.
The rank assertion follows either by the given triangular
flag or by reduction and Nakayama over the local Artin
ring.  In any flag with the prescribed ordered
characters, the $i$th flag submodule is exactly the sum
of the first $i$ of these eigenspaces: the idempotents
preserve the flag and have the indicated ranks on its
successive quotients.  This proves uniqueness, and also
descends the flag through any finite coefficient
extension used to apply the theorem, since the
projectors (G57.3) are already defined over $C$.

For $g\in G_{\mathbb Q_\ell}$, translate the flag
by $r_C(g)$.  Since $L/\mathbb Q_\ell$ is Galois,
the new flag is still $G_L$-stable.  On $I_L$ its
ordered characters are unchanged: cyclotomic
characters are invariant under conjugation by $g$.
Uniqueness makes the translated flag equal to the
original flag.  It is therefore stable under the
whole rational local group.

Write its ordered rank-one quotient characters as
$a_0,a_1,a_2$.  Put $h=(e,0,-e)$ and
$t_i=a_i\kappa_\ell^{-h_i}$ on inertia.
Each $t_i$ is trivial on $I_L$, so it is a character
of the finite group $I_\ell/I_L$ with values in
$C^\times$.  The reduction of the flag is the
standard flag of the seed, by the same uniqueness
argument.  Hence the reduction of $t_i$ is respectively
$\psi_I,1,\psi_I^{-1}$.
Divide $t_i$ by this constant coefficient character.
Its image lies in $1+\mathfrak m_C$ and is finite.
This group has no nontrivial finite-order elements:
the finite nilpotent logarithm is a group isomorphism
to the additive group $\mathfrak m_C$, which has
characteristic zero and is torsion free as an additive
group.  The quotient characters are therefore exactly
the constants just specified.  This proves (G57.1).
\end{proof}

\begin{theorem}[AC-G57: every split-component point has ordered ordinary inertia]
\label{AC-G57}
Put $R_*=U/P_*$ and let $r$ be its universal local
representation.  For every $s\in I_\ell$ and every
triple $s_0,s_1,s_2\in I_\ell$ one has in $R_*$
\[
\begin{split}
 \det(X-r(s))&=\prod_{i=0}^2(X-\gamma_i(s)),\\
 (r(s_0)-\gamma_0(s_0)I)
 (r(s_1)-\gamma_1(s_1)I)
 (r(s_2)-\gamma_2(s_2)I)&=0,
\end{split}                                                    \tag{G57.4}
\]
where $(\gamma_0,\gamma_1,\gamma_2)=(\gamma,1,\gamma^{-1})$.
Consequently every characteristic-zero field-valued
$O$-algebra point $f:R_*\to F$ has a unique ordered
ordinary flag: the full local group preserves the
flag, and its inertia characters in order are
$\gamma,1,\gamma^{-1}$.
In particular this holds for every actual local
point used in AC-G54-model.
\end{theorem}
\begin{proof}
For each $N\geq1$, compose the actual embedding
$\Phi$ with reduction to
$C_N=E[[T_1,\ldots,T_{12}]]/(T_1,\ldots,T_{12})^N$.
These are the continuous Artin points at the actual
seed justified in AC-G54.  They kill the fixed-type
ideal and are potentially crystalline of its type.
Lemma AC-G57-Artin applies to each one.
Upper triangular matrices with the specified
ordered inertia characters have the displayed
characteristic polynomial.  For the product identity,
the rightmost factor sends the whole module into the
second flag submodule, the middle factor sends that
submodule into the first, and the leftmost factor
kills the first.  Thus both identities in (G57.4)
hold modulo every power of $(T_1,\ldots,T_{12})$.
The formal-series ring is separated, so they hold
in that ring; the injectivity of $\Phi$ gives them
in $R_*$.  This proves integral polynomial identities,
without a density assertion about characteristic-zero
closed points and without identifying components away
from the formal germ.

We give the field-valued linear algebra in full.
The structure map $O\to F$ is injective: every
nonzero element of $O$ is a unit times a power of
a uniformizer, and a positive power of that
uniformizer is a unit times $\ell$; characteristic
zero prevents its image from vanishing.
Choose $s\in I_\ell\cap\ker\psi$ with the
three values $d_i=\gamma_i(s)$ distinct.
By the first identity in (G57.4), choose a basis in
which $f(r(s))=D=\operatorname{diag}(d_0,d_1,d_2)$.
For $x\in I_\ell$, put $N=f(r(x))$ in this basis
and $p_i=\gamma_i(x)$.
Apply the second identity successively to
$(x,s,s)$, $(s,x,s)$ and $(s,s,x)$ to get
\[
\begin{split}
 (N-p_0I)(D-d_1I)(D-d_2I)&=0,\\
 (D-d_0I)(N-p_1I)(D-d_2I)&=0,\\
 (D-d_0I)(D-d_1I)(N-p_2I)&=0.
\end{split}\tag{G57.5}
\]
In the first equation the right diagonal product has
only its zeroth diagonal entry nonzero; thus
$N_{00}=p_0$ and $N_{10}=N_{20}=0$.
In the third equation the left product has only
its second entry nonzero; thus $N_{21}=0$ and
$N_{22}=p_2$.  The $(1,1)$ entry of the middle
equation gives $N_{11}=p_1$.
All cancelled factors are products of distinct
$d_i-d_j$.  Every inertia matrix is therefore
upper triangular with the required diagonal.

For arbitrary $g\in G_{\mathbb Q_\ell}$, set
$A=f(r(gsg^{-1}))$ and $k=f(r(g))$ in this basis.
The matrix $A$ is upper triangular with diagonal
$d$, because inertia is normal and each $\gamma_i$
is a character of the full local group.
The representation identity gives $Ak=kD$.
Its $(2,0)$ entry first gives
$(d_2-d_0)k_{20}=0$; its $(2,1)$ entry then gives
$(d_2-d_1)k_{21}=0$, and its $(1,0)$ entry,
using $k_{20}=0$, gives
$(d_1-d_0)k_{10}=0$.
Thus $k$ is upper triangular.  The full local
group preserves the ordered flag.  Any other flag
with the same ordered inertia characters is the
ordered sum of the eigenspaces of $r(s)$, hence is
equal to this flag.  This finishes the proof.

The last calculation is precisely the matrix part
of the source
\path{Development/AutomorphicArithmetic/AutomorphicVarieties/PadicHodgeTheory/CommutingUpperTriangularDistinct.lean},
\texttt{lemma2b}.  Here its polynomial hypotheses
have been proved for the new ideal $P_*$ directly;
the old \texttt{D.On U r} premise is not silently
transferred to a different local ring.
\end{proof}

\begin{corollary}[AC-G57-TypeII: the ordinary condition for point transfer]
\label{AC-G57-TypeII}
At the newly constructed split primes, the actual
ordered ordinary hypothesis in AC-G48-TypeII and
AC-K29-point is now supplied for every actual
cubic Type II point, retaining the middle-character,
endpoint-determinant, residual-kernel and published
source hypotheses of those results.
The original selected-component factorization,
newform and level conclusion therefore extend to
these split primes.
\end{corollary}
\begin{proof}
The actual local point kills $P_*$, so AC-G57 gives
the ordered flag.  For a block Type II point, its
middle line has inertia character one; its endpoint
ordinary line and quotient have characters
$\gamma$ and $\gamma^{-1}$.  This follows either
by intersecting the ordered eigenspace sums with
the endpoint block, or by the three distinct
inertia eigenvalues used in the proof above.
Untwist the endpoint block by $\xi^{-1}$ as in
AC-K29-kill.  The character $\xi$ is unramified,
so this operation preserves the ordered inertia
characters.

We check crystallinity on the specific source
subgroup, rather than merely on some finite
subgroup.  In the literal source
\path{PadicHodgeTheory/IsInertiaInfiniteCyclotomic.lean},
$H_L=\ker\eta\cap\ker\tau$ and
$\Psi_I=\kappa_\ell^e\tau_I$.
Here $\eta=1$ locally, and the fixed inertial
type is $\tau_I\oplus1\oplus\tau_I^{-1}$.
It is trivial on $I_\ell\cap H_L$.
Let $F$ be the local field with group $H_L$,
and choose a finite Galois extension $L/F$ over
which the actual specialized representation is
crystalline.  Its full crystalline period module
$D$ is free of rank three over $L_0\otimes E'$,
where $E'$ is a finite field of coefficients
and $L_0$ is the maximal unramified subfield.
The action of $\operatorname{Gal}(L/F)$ is
semilinear on the $L_0$ factor.  Its inertia
subgroup acts trivially on $L_0$ and, by the
specified trivial inertial type, trivially on
$D$.  The remaining unramified Galois action
therefore descends $D$ to a free rank-three
module over $F_0\otimes E'$ by ordinary
finite Galois descent, componentwise if that
coefficient algebra is a product of fields.
These invariants are exactly
$(B_{\mathrm{cris}}\otimes r)^{G_F}$.
They have the full rank, so the representation
is crystalline on $H_L$.
This is the period-module definition of the
fixed potentially crystalline inertial type;
there is no semistable monodromy term because
the representation is potentially crystalline.
The unramified endpoint untwist is crystalline
and preserves this conclusion.

The original residual congruence,
normalization and actual inertia-kernel killing
are supplied by AC-K29-kill.  Under the retained
determinant identity the untwisted endpoint block
has determinant one, as required by AC-G48-TypeII.
That proposition proves membership in the original
selected local prime.  The finite coefficient-ring
construction and actual universal factorization
in AC-K29-point now apply verbatim to this verified
point, and AC-K29-typeII gives the retained
newform and level conclusion.
\end{proof}

\paragraph{Scope and geometric dictionary.}
The all-point ordered local condition missing in
AC-G54 is now proved.  This does not prove M2,
Type I automorphy, or vanishing of the projected
mixed tangent space.  The Euler conclusions retain
their own ordinary good-level, cubic-distinction
and local-factor hypotheses.  The proof passes
from infinitesimal neighborhoods of the actual
crystalline component to identities in its
coordinate ring using an injective formal chart.
The subsequent flag is obtained by explicit
spectral idempotents and matrix identities.
No new Grothendieck archival identification or
geometric comparison with the adele-class space
is asserted.
